\chapter{Lexer}
\label{ch:lexer}

% Source: src/zlex.c, include/ztok.h

\section{Overview}

The lexer (\texttt{zlex.c}) converts the raw source text into a flat list of \texttt{ZToken} objects consumed by the parser.
It is invoked once per source file and operates in a single forward pass.

It returns a \texttt{ZTokenStream}  -  a cursor wrapper around the token list that the parser can advance through.

\section{ZLexer State}

Internally the lexer holds a \texttt{ZLexer} struct:

\begin{verbatim}
typedef struct {
    ZState  *state;
    char    *program;   // full file buffer
    char    *current;   // current scan position
    ZToken **tokens;    // dynamic array of produced tokens
    char    *line;      // start of current line (for error reporting)
    usize    row, col;  // current position
    bool     sawNewline; // true if a newline was seen since last token
} ZLexer;
\end{verbatim}

\section{Token Types}

All token types are defined in \texttt{include/ztok.h} with the X-macro pattern and used in \texttt{include/zinc.h} to define the \texttt{ZTokenType} enum.
They fall into several categories:

\begin{itemize}
  \item \textbf{Keywords}: \zn{if}, \zn{else}, \zn{for}, \zn{return}, \zn{struct}, \zn{pub}, \zn{foreign}, \zn{use}, \zn{match}, \zn{defer}, \zn{goto}, \zn{macro}, \zn{type}, \zn{interface}, \zn{with}, \zn{and}, \zn{or}, \zn{not}
  \item \textbf{Literals}: integer, float, string, boolean (\zn{true}/\zn{false}), \zn{none}
  \item \textbf{Identifiers}
  \item \textbf{Operators \& punctuation}: \zn{+}, \zn{-}, \zn{*}, \zn{/}, \zn{:=}, \zn{->}, \zn{::}, \zn{==}, \zn{!=}, \zn{<=}, \zn{>=}, \ldots
\end{itemize}

\section{Keyword Hashing}

To avoid a linear scan through the keyword list on every identifier, the lexer uses a fixed-size hash map:

\begin{verbatim}
#define HASHMAP_TOK_LEN  256
#define HASHMAP_TOK_MASK (HASHMAP_TOK_LEN - 1)

typedef struct { const char *keyword; ZTokenType type; } KeywordEntry;
static KeywordEntry keywordEntries[HASHMAP_TOK_LEN];
\end{verbatim}

Keys are hashed with FNV-1a (32-bit). Collisions are resolved by open addressing with linear probing.
The table is populated once at startup via \texttt{initKeywords()}.
In standard implementation of an hashmap the lenght cannot be fixed and the entire hashmap must be grow. For a lexer, where symbols and keywords are fixed, it can be fixed to minimize implementation error.

\section{Newline Sensitivity}

Zinc uses newlines as implicit statement terminators in certain contexts to avoid typing ';' at every statement.
The \texttt{sawNewline} flag tracks whether at least one newline was consumed between the previous token and the current one.
The parser consults this flag when deciding whether to insert an implicit semicolon. The missing semicolon makes the parser harder because it doesn't know if it can continue and try to parse the next token or the statement is terminated. In some cases (like for \zn{for} and \zn{if} ...) is simpler because the terminator is the '\}'.

Here is an example of a 'difficult' program to parse without an explicit semicolon:

\begin{verbatim}
val := 10
*MyStruct obj = none
\end{verbatim}

When the parser reaches the token 10 it sees * and it thinks that it is parsing a multiplication.

\section{Error Reporting}

Lexical errors (unterminated strings, unknown characters) are reported through \texttt{ZState} via the shared \texttt{error()} function.
The lexer records the \texttt{row}/\texttt{col} of every token so that error messages can point to the exact source location.
